\documentclass{IEEEtran}
\usepackage{fancyvrb}
\fvset{numbers=left,fontsize=\fontsize{8}{10}\selectfont,numbersep=6pt}

\usepackage{amsmath}
\usepackage{float}
\usepackage{hyperref}

\newcommand{\pannelwidth}{0.32}
\setlength{\parskip}{1ex}
\setlength{\baselineskip}{0.1ex}

\title{Title of your reflective report}
%\author{Your Name}
\author{Your Role Number}

\date{\today}

\begin{document}

\maketitle
\begin{abstract}
    Your abstract summarizing your reflective report. Reflective reports are not anonymized, do not use your roll number.
\end{abstract}

\section{Introduction}

\IEEEPARstart{Y}{our} introduction goes here. Simply start writing your document and use the Recompile button to view the updated PDF preview. Examples of commonly used commands and features are listed below, to help you get started.

Introduction to the question should show reasoning why you chose to implement this algorithm(s). Explain how your algorithm(s) could be used to address what biomedical questions and the significance of your work. 

Once you're familiar with the editor, you can find various project setting in the Overleaf menu, accessed via the button in the very top left of the editor. To view tutorials, user guides, and further documentation, you can visit this \href{https://www.overleaf.com/learn}{help library}.

\section{Methods and Results}

This section includes how the problem has been analyzed. What algorithm designing approaches you applied to address the question. You may include a brief description of your algorithm(s) design, a flow chart to show the input/output as well as step by step design of your algorithm(s). This part may include the source data, publication or data link including database ID for the dataset used. If you use simulated data, you need to clearly state how you simulated the data. Here, you must describe how you benchmarked your algorithm.

\section{Reflection}

Reflection on your work should be both as an individual and as a group. What did you and your group did well? What and how could be improved?

\section{Discussion}

This section should include discussion about your algorithm(s). Some aspects you can discuss include: how your algorithm(s) could be improved? Maybe it is only working in small scale data and cannot scale up? Any potential ideas to optimize this? Is there other published or publicly available algorithms (*e.g.* on github), addressing the same question? What are the advantages or disadvantages of your algorithm(s) compared to others? If your algorithm(s) is novel and better than others, you can translate this course project to a publication.

\section{References}

Add your used references. Any (consistent) format is fine.

The total length of your report cannot exceed 4 pages.

\end{document}